\uuid{tEBA}
\titre{Changement de variables}
\auteur{hanani ; mzari}
\organisation{exo7}
\licence{CC BY-NC-SA 4.0}
\datecreate{2020-01-20}
\niveau{L1}
\module{Analyse}
\chapitre{Calcul d'intégrales}
\sousChapitre{Primitives diverses}
\difficulte{3}
\theme{Primitives, changement de variables}
\exo7qcmid{lille-2.primitives.046}

\contenu{
\begin{question}
\qclassification{127.11}
\qtags{motcle=changement de variables}
Le changement de variable $t=\sin x$ donne :
\begin{answers}
\bad{$\displaystyle \int \frac{\cos ^3x\mathrm{d}x}{1+\sin ^2x}=\cos ^3 x\int \frac{\mathrm{d}t}{1+t^2}$.}
\bad{$\displaystyle \int \frac{\cos ^3x\mathrm{d}x}{1+\sin ^2x}=\cos ^3 x.\arctan (\sin x)+k$, $k\in \Rr$.}
\good{$\displaystyle \int \frac{\cos ^3x\mathrm{d}x}{1+\sin ^2x}=\int \frac{1-t^2}{1+t^2}\mathrm{d}t$.}
\good{$\displaystyle \int \frac{\cos ^3x\mathrm{d}x}{1+\sin ^2x}=-t+2\arctan t+k$, $k\in \Rr$.}
\end{answers}
\vskip2mm
\begin{explanations}
D'abord, $\displaystyle \frac{\cos ^3x}{1+\sin ^2x}=\frac{\cos x(1-\sin ^2x)}{1+\sin ^2x}$. Ensuite $t=\sin x\Rightarrow\mathrm{d}t=\cos x\mathrm{d}x$ et
$$\int \frac{\cos ^3x\mathrm{d}x}{1+\sin ^2x}=\int \frac{1-t^2}{1+t^2}\mathrm{d}t=\int \left(-1+\frac{2}{1+t^2}\right)\mathrm{d}t=-t+2\arctan t+k,\; k\in \Rr.$$
\end{explanations}
\end{question}
}
